\chapter{机器学习究竟在学习什么}

\begin{quote}\small
\textit{本章摘要。}本章定义样本、任务、模型、损失和风险，说明不同损失对应不同预测对象，并讨论正则化、可识别性及预测到决策的边界。
\end{quote}

上一章给出了向量、导数和概率的语言，但能计算一个梯度并不意味着已经定义了值得学习的问题。现在需要把设备观测、未来事件和维护动作分开：观测决定模型能看见什么，标签决定它试图回答什么，损失决定怎样的错误更值得避免。三者不一致时，即使训练完全收敛，也可能得到一个回答错误问题的模型。

\paragraph{本章导语。}
同一段振动记录，可以用来判断轴承是否故障，也可以用来预测下一时刻的振幅。这两个问题需要不同的标签，算错后的后果也不同。本章先把这些差别写清楚，再说明模型中的参数如何通过损失函数学出来。随后讨论正则化、可识别性和部署反馈：为什么训练数据上看起来很好，换一台设备却可能失效？后文介绍各种模型时，都要回到输入、目标和错误代价这几件事。

\section{学习对象：数据、任务与模型}

% auxiliary-resource-guide
本章末的“辅助学习材料”列出与核心方法对应的讲解和实践入口，可在阅读相关推导后，按所列小节对照学习。

\paragraph{从数据到任务。}
数据（data）是对对象、过程或环境的观测记录。一个样本（sample）是一次相对完整的观测，通常记为$\bm{x}\in\R^d$；监督学习中还伴随标签$y$，于是样本写为$(\bm{x},y)$。这里的“样本”不是数据表的一行这么简单：一次设备运行周期、一张图像、一个句子或一段传感器窗口都可以成为样本，关键在于它是否对应一个明确的学习单位。

任务由输入空间$\mathcal{X}$、输出空间$\mathcal{Y}$以及评价准则共同决定。分类的输出是离散类别，回归的输出是连续数值，预测性维护可能输出剩余寿命，异常检测则可能没有逐样本标签。相同的数据可以支持不同任务。例如，振动信号既可以用于故障分类，也可以用于生成未来波形或估计设备状态。

\paragraph{模型、参数与预测。}
模型（model）是带有可调参数的函数族
\begin{equation}
 f_{\bm{\theta}}:\mathcal{X}\rightarrow\mathcal{Y},
 \qquad \bm{\theta}\in\Theta.
\end{equation}
给定输入$\bm{x}$，确定性模型直接输出$\hat{y}=f_{\bm{\theta}}(\bm{x})$。参数$\bm{\theta}$不是由研究者凭空指定的答案，而是由训练数据和训练目标共同决定的可学习量。模型结构规定了“允许怎样的映射”，参数则规定了“当前采用哪一个映射”。

概率模型输出条件分布$p_{\bm{\theta}}(y\mid\bm{x})$。分类时可以把它理解为各类别的置信度，回归时则可以同时表示预测均值和噪声大小。确定性输出可以看作概率模型的特殊情形，例如$p(y\mid\bm{x})$集中在单点$f_{\bm{\theta}}(\bm{x})$上。因此，确定性模型与概率模型不是互斥的两个世界，而是表达不确定性程度不同的两种建模方式。

\section{训练目标：从预测到可优化问题}
有限样本上的拟合与总体上的预测能力需要分开讨论，经验分布到真实分布的一致收敛是统计学习理论的基础问题之一\cite{vapnik1971uniform}。神经网络的反向传播则提供了多层可微模型求梯度的有效途径\cite{rumelhart1986backprop}；求得梯度并不自动保证泛化。
\paragraph{损失函数与训练目标。}
损失函数（loss function）$\ell(\hat{y},y)$衡量一次预测的代价。训练集$\mathcal{D}=\{(\bm{x}_i,y_i)\}_{i=1}^n$上的经验风险为
\begin{equation}
 \hat{R}(\bm{\theta})=\frac{1}{n}\sum_{i=1}^{n}\ell\bigl(f_{\bm{\theta}}(\bm{x}_i),y_i\bigr).\label{eq:foundations-empirical-risk}
\end{equation}
以式\eqref{eq:foundations-empirical-risk}为依据，最基本的训练目标是$\argmin_{\bm{\theta}}\hat{R}(\bm{\theta})$。回归常用平方损失$\frac12(\hat{y}-y)^2$，分类常用交叉熵
\begin{equation}
 \ell_{\mathrm{CE}}=-\sum_{k=1}^{K}\mathbb{1}[y=k]\log p_{\bm{\theta}}(y=k\mid\bm{x}).
\end{equation}
损失函数不是“模型好坏”的全部定义。它还必须与实际错误代价、标签质量和部署决策相匹配。工业场景中漏报和误报往往代价不同，单纯追求平均准确率可能掩盖真正重要的风险。

\subsection{从总体风险到有限样本目标}
设$(X,Y)\sim P$描述部署环境中的输入与标签，$X\in\R^d$，$Y$可为实数或类别。总体风险定义为
\begin{equation}
 R(f)=\E_{(X,Y)\sim P}[\ell(f(X),Y)].
\end{equation}
这里的期望包含尚未观察到的未来样本，因此通常不能直接计算。若训练样本独立同分布于$P$，对一个事先固定且损失可积的$f$，有$\E[\hat R(f)]=R(f)$。若损失方差有限，独立性进一步给出
\begin{equation}
 \Var(\hat R(f))=\frac{1}{n^2}\sum_{i=1}^n\Var(\ell(f(X_i),Y_i))
 =\frac{\Var(\ell(f(X),Y))}{n}.
\end{equation}
这解释了平均损失为什么可用于估计风险，也暴露了限制：相邻设备窗口相关时有额外协方差项；$f$由同一训练集选择后，上述固定模型的无偏论证不能直接用于其训练误差。验证与测试的必要性由此产生，而不是来自惯例。

\subsection{损失决定预测对象：均值、中位数与概率}
平方损失的条件风险可分解为$\tfrac12\Var(Y\mid X=x)+\tfrac12(a-\E[Y\mid X=x])^2$，所以最优点预测是条件均值。若使用绝对值损失，令$F_x(a)=P(Y\leq a\mid X=x)$；在条件分布连续且期望有限时，
\begin{equation}
 \frac{d}{da}\E[|Y-a|\mid X=x]
 =P(Y<a\mid x)-P(Y>a\mid x)=2F_x(a)-1.\label{eq:foundations-absolute-risk-gradient}
\end{equation}
令式\eqref{eq:foundations-absolute-risk-gradient}为零，最优点满足$F_x(a)=1/2$。有原子概率时应使用左右导数，最优点满足$P(Y<a\mid x)\leq1/2\leq P(Y\leq a\mid x)$，可能不唯一。

更一般地，令$0<\tau<1$，分位数损失为$\rho_\tau(u)=u(\tau-\mathbb1[u<0])$，其中$u=Y-a$。预测偏低时斜率代价为$\tau$，偏高时为$1-\tau$，条件风险对$a$的导数为$F_x(a)-\tau$，于是最优解是条件$\tau$分位点。若维护决策特别关心寿命偏高估计，可通过适当分位点表达这种不对称性，但分位点的选择仍须来自真实决策成本。

对$K$类问题，记真实条件概率为$p_k=P(Y=k\mid X=x)$，预测概率为$q_k>0$、$\sum_kq_k=1$。条件交叉熵为
\begin{equation}
 \E[-\log q_Y\mid X=x]=-\sum_kp_k\log q_k
 =H(p)+\KL(p\|q).
\end{equation}
因此在可表达真实分布且达到总体最优的理想条件下，交叉熵学习整个条件分布；只以准确率训练则只需找最大概率类别。有限数据下输出在$[0,1]$内不等于已经校准。

\section{训练目标与结构约束}
\paragraph{正则化与归纳偏置。}
当模型足够灵活时，它可能记住训练样本而没有学到可迁移的规律。正则化（regularization）通过在经验风险之外加入偏好来限制解，例如
\begin{equation}
 \min_{\bm{\theta}}\;\hat{R}(\bm{\theta})+\lambda\Omega(\bm{\theta}),
 \qquad \Omega(\bm{\theta})=\frac12\|\bm{\theta}\|_2^2.\label{eq:foundations-regularized-risk}
\end{equation}
式\eqref{eq:foundations-regularized-risk}中，$\lambda$控制拟合数据和保持模型简单之间的权衡。权重衰减、数据增强、早停、Dropout和共享参数都可以被理解为不同形式的归纳偏置（inductive bias）：它们把某些解视为更合理，从而在有限数据下提高泛化可能性。

\begin{example}
 对于温度传感器预测下一时刻温度，模型并不需要学习所有历史序列到未来数值的任意映射。若已知相邻时刻通常连续变化，可以采用平滑性偏置；若设备存在周期，则可以把周期特征加入输入。模型的“聪明”来自数据，也来自这些合理的限制。
\end{example}

\paragraph{惩罚形式与约束形式如何联系。}
也可把偏好写成$\min_\theta\hat R(\theta)$且$\Omega(\theta)\leq c$。其拉格朗日函数为$L(\theta,\lambda)=\hat R(\theta)+\lambda(\Omega(\theta)-c)$，$\lambda\geq0$。固定$\lambda$后，常数$-\lambda c$不影响参数最优解，便得到惩罚目标。若目标与约束凸并满足适当的严格可行性条件，最优解可满足
\begin{equation}
 0\in\partial\hat R(\hat\theta)+\lambda\partial\Omega(\hat\theta),\qquad
 \Omega(\hat\theta)\leq c,\qquad
 \lambda(\Omega(\hat\theta)-c)=0.
\end{equation}
$\partial$是允许不可微凸函数的次梯度集合。互补条件表示：限制未起作用时乘子可为零；限制起作用时，数据拟合的改善必须补偿复杂度增加的代价。非凸网络中不能据此声称任意约束问题都与某个惩罚问题全局等价；$\lambda$与$c$也不必一一对应。

以一维模型、$\hat R(w)=\tfrac12(w-a)^2$为例，$a$是无正则化估计。加入$\lambda w^2/2$后，导数为$(w-a)+\lambda w$，所以$\hat w=a/(1+\lambda)$。这展示了收缩的具体含义：不是自动删掉噪声，而是把所有估计向零拉近。若真实系数远离零，收缩会引入偏差；只有减少的不稳定性足够大时，总预测风险才可能下降。

\paragraph{从训练到部署及反馈。}
沿图\ref{fig:foundations-learning-loop}从采集到上线的路径，可以依次看四个阶段。先把设备记录整理成输入和标签，再选择能处理这些输入的模型，通过训练求出参数，最后用这些参数预测新记录。上线后的错误又会提示我们补数据、改标签或重新训练。评价时除了训练损失，还要检查新工况下的表现、概率是否校准、计算是否及时，以及误报和漏报的后果。后续章节将分别展开这些步骤。

\begin{figure}[htbp]
\centering
\begin{tikzpicture}[>=Latex, node distance=0.85cm, every node/.style={font=\small}]
 \node[draw, rounded corners, fill=blue!10, minimum width=2.15cm, minimum height=0.9cm] (data) {\shortstack{数据\\$\mathcal D$}};
 \node[draw, rounded corners, fill=orange!12, right=of data, minimum width=2.35cm, minimum height=0.9cm] (task) {\shortstack{任务定义\\$X,Y,\ell$}};
 \node[draw, rounded corners, fill=green!12, right=of task, minimum width=2.35cm, minimum height=0.9cm] (model) {\shortstack{模型\\$f_\theta$}};
 \node[draw, rounded corners, fill=yellow!20, right=of model, minimum width=2.35cm, minimum height=0.9cm] (train) {\shortstack{训练\\$\argmin_\theta$}};
 \node[draw, rounded corners, fill=red!10, right=of train, minimum width=2.15cm, minimum height=0.9cm] (deploy) {\shortstack{推理\\与部署}};
 \draw[->, thick] (data)--(task); \draw[->, thick] (task)--(model); \draw[->, thick] (model)--(train); \draw[->, thick] (train)--(deploy);
 \draw[->, thick, dashed] (deploy.north) .. controls +(0,0.72) and +(0,0.72) .. node[above,font=\scriptsize]{反馈、漂移、错误分析} (data.north);
\end{tikzpicture}
\caption{工业数据智能的训练、部署与反馈：部署反馈会重新影响数据和任务定义。}
\label{fig:foundations-learning-loop}
\end{figure}

\paragraph{学习对象与可识别性。}
同一组观测可以对应不同的学习对象。若输入是温度窗口，目标可以是下一时刻温度、未来一小时最大温度、设备健康分数或故障概率。它们需要不同标签、不同时间边界和不同损失函数。若标签由未来事件定义，训练样本必须避免使用事件发生之后才出现的信息。

还需区分“可预测”和“可识别”。一个变量可能与故障高度相关，却只是设备编号的替代变量；当设备更换后，相关性就会消失。好的特征不是在固定数据表上有高相关系数，而是在数据生成机制改变后仍保留任务所需的信息。

\paragraph{确定性预测与概率预测。}
平方损失学习条件均值，绝对值损失学习条件中位数，分位数损失学习指定分位点。这说明损失函数实际上规定了预测对象。若同一输入对应多个合理输出，点预测会把它们平均在一起；概率预测则尝试给出完整条件分布。部署时若动作成本不对称，还应从预测分布计算决策风险，而不是直接选择最大概率类别。

\paragraph{可识别性不是相关性的另一个名字。}
严格地说，参数可识别要求不同参数对应不同的可观测分布；仅讨论确定性预测时，可要求$f_\theta(X)=f_{\theta'}(X)$几乎处处成立只能推出$\theta=\theta'$。例如若所有样本均满足$X_2=2X_1$，则
\begin{equation}
 w_1X_1+w_2X_2=(w_1+2w_2)X_1.
\end{equation}
数据只能确定组合$w_1+2w_2$，无法分别确定两个系数。正则化可以选出某个解，却没有创造额外观测证据。参数不可识别、跨设备相关性失效以及因果关系不可识别是不同问题；它们共同说明准确预测并不自动赋予参数机制解释。

\paragraph{从条件概率推导动作。}
令动作$a\in\{0,1\}$分别表示不告警与告警，正确动作成本暂取零，误报与漏报成本分别为$c_{FP},c_{FN}>0$。给定校准到目标环境的故障概率$p$，条件风险为
\begin{equation}
 r(1\mid x)=c_{FP}(1-p),\qquad r(0\mid x)=c_{FN}p.\label{eq:foundations-action-risk}
\end{equation}
比较式\eqref{eq:foundations-action-risk}中的两种风险得到告警条件$p\geq c_{FP}/(c_{FP}+c_{FN})$，等号处两种动作成本相同。若漏报成本是误报的九倍，阈值为$0.1$而不是$0.5$。此推导假设动作成本固定且概率对应部署人群；若告警后还会改变故障概率，或维护容量有限，就需要扩展动作与后果模型，不能仅调整一个阈值解决全部问题。

\paragraph{模型选择的四个约束。}
模型选择至少同时受四类约束：数据规模与标注质量、任务中的依赖结构、训练与推理资源、错误的业务代价。复杂模型并不自动优于简单模型；它只有在表示了简单模型无法表示的结构时才有必要。后文讨论的序列、视觉和语言模型，实际上是在不同结构约束下对这四个问题的回答。

本章的风险、损失与正则化可结合下面两份非论文材料复习：前者帮助把目标写成可计算的学习问题，后者把模型训练与泛化评价连接起来。
\section*{辅助学习材料}
\begingroup
\emergencystretch=3em
\begin{itemize}
\item \textbf{Linear Models}；作者/平台：scikit-learn Developers；英文；官方文档。重点阅读 Ordinary Least Squares、Ridge、Lasso 与 Logistic Regression 小节，配合本章的线性预测、正则化和分类损失理解参数约束与风险最小化。\href{https://scikit-learn.org/stable/modules/linear\_model.html}{在线阅读}。
\item \textbf{Cross-validation: evaluating estimator performance}；作者/平台：scikit-learn Developers；英文；官方文档。重点阅读训练/验证/测试划分、交叉验证和数据泄漏说明，帮助把本章的经验风险、部署风险与可复现实验边界落实为流程。\href{https://scikit-learn.org/stable/modules/cross\_validation.html}{在线阅读}。
\item \textbf{损失函数（Loss Function）}；知乎；中文解说。比较平方损失与交叉熵，核对本章预测对象；标题与主题据必应索引，原页受限。\href{https://zhuanlan.zhihu.com/p/261059231}{在线阅读}。
\item \textbf{损失函数（lossfunction）的全面介绍（简单易懂版）}；CSDN；中文博客。阅读 L1、MSE 与交叉熵说明，比较不同误差的代价；已核对公开页面标题与索引摘要。\href{https://blog.csdn.net/weixin\_57643648/article/details/122704657}{在线阅读}。
\item \textbf{2021 - 类神经网络训练不起来怎么办(四) 损失函数 (Loss) 也可能有影响}；李宏毅；bilibili；中文课程。选看第22集，对照本章损失与学习目标。2026年10月4日官方公开元数据显示整套课程累计播放约299.5万，并非本分集播放量；已核对目录与计数，未逐帧观看。\href{https://www.bilibili.com/video/BV1Wv411h7kN/?p=22}{观看分集}。
\end{itemize}
\endgroup
\section{本章小结：学习的对象由规则决定}
机器学习并非从数据中提取一个与任务无关的唯一答案。输入和标签定义可观察的问题，模型族限定可表达的答案，损失确定希望估计的统计对象，正则化与优化再从有限数据允许的解中作出选择。总体风险与经验风险、预测与解释、概率与动作必须分开，才能判断一个低损失结果到底说明了什么。

下一章用可以手算的经典模型练习这些概念。平方损失怎样导出正规方程，惩罚为什么使系数收缩或变为零，逻辑回归的梯度怎样从交叉熵得到，都将逐步展开。先把简单模型算明白，后面换用复杂模型时，才能说清它解决了哪一种原先解决不了的问题。
